\section*{Chapter \ref{ch:b3:factor-endowments} --- Factor Endowments}

\begin{solution}{exo:b3:factor-endowments:1}
China (19.1~\% of capital against 17.5~\% of output), Japan (4.4 against 3.8),
Russia (5.8 against 3.4), France (3.0 against 2.4), Indonesia (3.2 against 2.4), the
United Kingdom (2.4 against 2.3), Italy (2.9 against 1.9) and South Korea (2.0
against 1.8). % figdata: endow
By HOV, with output standing in for spending, they should be net exporters of
capital services. The United States, India, Germany and Brazil have smaller
shares of capital than of output.
\end{solution}

\begin{solution}{exo:b3:factor-endowments:2}
Steel uses $4/20 = \$0.2$ million of capital per worker, clothing $1/40 =
\$0.025$ million: steel is capital-intensive. If wages double, firms in both
industries substitute capital for labour; the ranking survives unless the
industries' abilities to substitute differ so much that their capital per
worker crosses, which the model rules out by assumption.
\end{solution}

\begin{solution}{exo:b3:factor-endowments:3}
$2w + r = 10$ and $w + 2r = 8$. Twice the first minus the second: $3w = 12$, $w = 4$;
then $r = 2$.
\end{solution}

\begin{solution}{exo:b3:factor-endowments:4}
In its two-good form, yes. But with many factors the comparison of capital per
worker in exports and imports is not the right test: the country could be
abundant in capital relative to land or skills even if it has little capital per
worker. Check the factor content of its trade against the HOV equation for every
factor, and check whether its technology differs from that of its partners.
\end{solution}

\begin{solution}{exo:b3:factor-endowments:5}
$0.7\hat w + 0.3\hat r = 10$ and $0.4\hat w + 0.6\hat r = 0$. From the second, $\hat r = -\tfrac23\hat
w$; then $0.7\hat w - 0.2\hat w = 10$, $\hat w = 20~\%$ and $\hat r = -13.3~\%$. So $\hat w = 20 >
\hat p_1 = 10 > \hat p_2 = 0 > \hat r = -13.3$: the \omterm{def:b3:factor-endowments:magnification}{magnification effect}.
\end{solution}

\begin{solution}{exo:b3:factor-endowments:6}
Labour: $2Q_1 + Q_2 = 100$; capital: $Q_1 + 2Q_2 = 80$. So $Q_1 = 40$, $Q_2 = 20$. With $K = 88$:
$Q_2 = 100 - 2Q_1$ and $Q_1 + 200 - 4Q_1 = 88$, so $Q_1 = 37.33$ and $Q_2 = 25.33$: good 1
falls by 6.7~\% and good 2, capital-intensive, rises by 26.7~\%, more than the 10~\%
rise in capital.
\end{solution}

\begin{solution}{exo:b3:factor-endowments:7}
Short run (\cref{prop:b3:factor-endowments:specific}): landowners gain in terms
of both goods; owners of capital in manufacturing lose in terms of both; workers'
wage rises by less than 20~\%, so they gain in manufactures and lose in grain.
Long run: capital moves into farming until its return is equal across sectors;
the conflict becomes one between factors, decided by which factor grain uses
intensively.
\end{solution}

\begin{solution}{exo:b3:factor-endowments:8}
The tested solution is \code{ex08_factor_content.py}, in the chapter's
\code{exercises/} folder. Capital: $4 \times 10 + 2 \times (-5) + 1 \times (-8) = 22$;
labour: $10 - 10 - 24 = -24$. % output: ex08_factor_content
The country exports the services of 22 units of capital, 5.5~\% of its endowment,
and imports those of 24 units of labour, 4.8~\% of its own: it trades as a
capital-abundant country.
\end{solution}

\begin{solution}{exo:b3:factor-endowments:9}
Wages for the same skill, in the same occupation, that differ persistently
between countries that trade freely with each other contradict it --- and they
differ many times over. They reject an assumption: identical technologies above
all (a worker with better technology earns more), together with the assumption
that both countries make the same goods.
\end{solution}

\begin{solution}{exo:b3:factor-endowments:10}
The system $a_{L1}w + a_{K1}r = p_1$, $a_{L2}w + a_{K2}r = p_2$ has determinant $D = a_{L1}a_{K2} -
a_{K1}a_{L2} \ne 0$, so by Cramer's rule exactly one solution: $w = (p_1a_{K2} - p_2a_{K1})/D$
and $r = (p_2a_{L1} - p_1a_{L2})/D$. It depends only on prices and technology, so two
countries with the same prices and technology, both making both goods, have the
same $w$ and $r$. Both are positive when $p_1/p_2$ lies strictly between
$a_{K1}/a_{K2}$ and $a_{L1}/a_{L2}$; otherwise one good cannot break even and the country
specialises.
\end{solution}

\begin{solution}{exo:b3:factor-endowments:11}
The tested solution is \code{ex11_hov_sign.py}, in the chapter's \code{exercises/}
folder. Exact HOV: signs match 1.00, slope 1.00. \omterm{def:b3:factor-endowments:factor-content}{Missing trade}: signs 1.00, slope
0.20. With noise: signs 0.57, slope 0.13. % output: ex11_hov_sign
A sign test cannot see \omterm{def:b3:factor-endowments:factor-content}{missing trade}, which shrinks contents without changing
their signs; the slope sees it at once. Noise of a size realistic for measured
data pushes the sign test toward one half.
\end{solution}

\begin{solution}{exo:b3:factor-endowments:12}
In a small open economy goods prices are set abroad and do not move; with both
goods still made, the zero-profit conditions fix $w$ and $r$
(\cref{def:b3:factor-endowments:fpe}). The inflow is absorbed by the \omterm{def:b3:factor-endowments:rybczynski}{Rybczynski
effect}: the labour-intensive sector expands and the other contracts. Conditions:
prices unaffected (a small economy, or immigrants who do not change world
prices), the same technology, and an inflow small enough that the economy still
makes both goods; a larger inflow would move it out of the range and wages
would fall.
\end{solution}

\begin{solution}{pb:b3:factor-endowments:boom}
\textbf{1.} Machines use $3/1 = 3$ units of capital per worker, textiles $1/2 = 0.5$:
machines are capital-intensive. % figdata: boom_params

\textbf{2.} Kestria $60/40 = 1.5$, Brenland $30/50 = 0.6$, the world $1{,}300/1{,}100 = 1.18$.
Kestria is capital-abundant, Brenland labour-abundant.

\textbf{3.} $3Q_M + Q_T = 60$ and $Q_M + 2Q_T = 40$: $Q_T = 60 - 3Q_M$, $Q_M + 120 - 6Q_M = 40$, so
$Q_M = 16$ and $Q_T = 12$ million. % figdata: boom_results

\textbf{4.} $p_M/p_T = 12/16 = 0.75$.

\textbf{5.} $3r + w = 0.75$ and $r + 2w = 1$: $r = 1 - 2w$, $3 - 5w = 0.75$, $w = 0.45$ and $r =
0.10$.

\textbf{6.} $0.75 \times 16 + 12 = 24$; $0.10 \times 60 + 0.45 \times 40 = 6 + 18 = 24$.

\textbf{7.} $3r + w = 4$ and $r + 2w = 3$: $w = 1$ and $r = 1$ kest.

\textbf{8.} $3Q_M + Q_T = 30$ and $Q_M + 2Q_T = 50$: $Q_M = 2$, $Q_T = 24$. Both countries
face the same prices with the same technology and make both goods, so the
zero-profit conditions give both the same factor prices.

\textbf{9.} GDP $= 4 \times 16 + 3 \times 12 = 100$; consumption $50/4 = 12.5$ machines and $50/3
= 16.67$ textiles; net exports $3.5$ machines and $-4.67$ textiles. $3.5 \times 4 = 14 =
4.67 \times 3$: trade balances.

\textbf{10.} Autarky: $w/p_T = 0.45$, $w/p_M = 0.45/0.75 = 0.60$; $r/p_T = 0.10$, $r/p_M =
0.13$. Trade: $w/p_T = 1/3 = 0.33$, $w/p_M = 1/4 = 0.25$; $r/p_T = 0.33$, $r/p_M = 0.25$.

\textbf{11.} $p_M/p_T$ rises from 0.75 to 1.33, by 77.8~\%. In textiles, the wage falls
from 0.45 to 0.33, by 25.9~\%, and the rental rises from 0.10 to 0.33, by 233~\%. So
$\hat r > \hat p_M > \hat p_T = 0 > \hat w$: the \omterm{def:b3:factor-endowments:magnification}{magnification effect}.

\textbf{12.} Owners of capital gain in terms of both goods; workers lose in terms
of both (\cref{prop:b3:factor-endowments:ss}).

\textbf{13.} $3Q_M + Q_T = 75$ and $Q_M + 2Q_T = 40$: $Q_M = 22$ and $Q_T = 9$ million.

\textbf{14.} Machines rise by $6/16 = 37.5~\%$, more than capital's 25~\%; textiles fall
by $3/12 = 25~\%$: \cref{prop:b3:factor-endowments:rybczynski}.

\textbf{15.} The new capital can only be used with workers; at unchanged factor
prices each machine needs one worker, and the workers come from textiles.

\textbf{16.} Prices are unchanged, so $w = r = 1$: workers' real wage does not change.

\textbf{17.} GDP $= 4 \times 22 + 3 \times 9 = 115 = 75 + 40$. Consumption $57.5/4 = 14.375$
machines and $57.5/3 = 19.17$ textiles; net exports 7.625 machines and $-10.17$
textiles.

\textbf{18.} Kestria's 75 million units are under 6~\% of the world's capital: its
outputs barely move world supply. If Kestria were large, the price of machines
would fall relative to textiles, and by Stolper--Samuelson its workers' real wage
would rise and the rental fall.

\textbf{19.} Capital: $3 \times 3.5 + 1 \times (-4.67) = 5.83$; labour: $3.5 + 2 \times (-4.67) =
-5.83$ (millions).

\textbf{20.} World GDP $= 1{,}300 + 1{,}100 = 2{,}400$, so $s = 100/2{,}400 = 0.0417$.

\textbf{21.} $60 - 0.0417 \times 1{,}300 = 60 - 54.17 = 5.83$ and $40 - 0.0417 \times 1{,}100 = 40 -
45.83 = -5.83$: equal to the factor content.

\textbf{22.} GDP 80, $s = 0.0333$; consumption 10 machines and 13.33 textiles; net
exports $-8$ and 10.67; factor content $3(-8) + 10.67 = -13.33$ of capital and $-8 +
21.33 = 13.33$ of labour; HOV: $30 - 43.33 = -13.33$ and $50 - 36.67 = 13.33$.

\textbf{23.} Measured factor contents much smaller than 5.83 million: Kestria
trading less than the model says. Causes: consumers who prefer home-made goods,
trade costs that keep goods at home, or a Kestrian technology that differs from
Brenland's, so that Kestria's own input requirements, not a common $A$, should be
used.

\textbf{24.} After the boom Kestria's textile output falls from 12 to 9 million
units, by 25~\%, while its workers' real wage is unchanged at world prices.
\end{solution}
